\section{Experimental design}
\label{sec:experiments}

All experiments are retrospective on the human-labeled dataset supplied for this project. No new
human-subject data are collected.

\subsection{E0: Data, splits, judges, and budgets}

\begin{itemize}
  \item \textbf{Dataset.} \todo{name, size, label structure (binary / graded / multi-annotator),
        availability of prompt and response texts, license}. E0 must be finished before any main
        experiment starts. Check the per-prompt forest condition of
        Section~\ref{sec:post-training} as part of E0.
  \item \textbf{Splits.} Split by prompt into an untouched test set $\mathcal T$ (20\%) and the pool
        (80\%, size $N$), before any acquisition, E2-driven choice or policy training. All
        estimation, acquisition and policy metrics are computed on $\mathcal T$. Every arm, including
        the all-human reference, trains on the pool only.
  \item \textbf{Judges.} J1 (primary) and J2 (secondary, E5 only) as defined in
        Section~\ref{sec:setup}. \todo{judge model(s); J1 needs token-level log-probabilities}.
  \item \textbf{Localization.} $Z_s=(p_J)$ and $Z_d=(p_J,t_J,o_J)$, stratified by task category.
        Pre-register the D1 degeneracy check using pool data only.
  \item \textbf{Budgets.} $n/N\in\{1,2,5,10,20\}\%$, adapted to $N$. For the estimation experiments,
        $50$ independent random reveals per budget.
\end{itemize}

\subsection{E1: Estimation of the human preference distribution}

\textbf{Methods.} Human-only (kernel on $Z_d$; logistic regression on $Z_d$); raw judge $\mJ$; scalar
recalibration of $p_J$ (Platt, isotonic); the linear map of \citet{vandenburg2025aligning} applied to
the full J1 output vector; panel calibration \citep{li2026calibrate} when several judges are available;
\sfsp; \dfsp.\\
\textbf{Metrics on $\mathcal T$.} Brier score, which equals $R_X$ up to a method-independent constant
(Lemma~\ref{lem:heldout}); NLL; RPS for graded labels; reliability diagrams and ECE; distance to the
empirical human distribution on multi-annotator items.\\
\textbf{Headline plot.} Test error versus the number of revealed human labels.

\subsection{\texorpdfstring{E2: Does the judge distribution carry information beyond $p_J$?}{E2: Does the judge distribution carry information beyond p\_J?}}

\begin{enumerate}[label=(\alph*)]
  \item \emph{Inference on $G$.} Using pool labels only, estimate the variable importance $G$ of
        $(t_J,o_J)$ given $p_J$ (Proposition~\ref{prop:scalar-vs-dist}). Use the cross-fitted,
        algorithm-agnostic estimator of \citet{williamson2023general}, including its sample-splitting
        version for the zero-importance null. The hypothesis and nuisance learners are fixed in
        advance. A fitted-score comparison is \emph{not} used as a test: its expectation is
        $G+R(\hmu_s)-R(\hmu_d)$.
  \item \emph{Calibration of the test.} A synthetic study keeps the real judge outputs but replaces
        the human labels with draws from $\Ber(g(p_J))$ (so $G=0$) and from $\Ber(g(p_J)+c\,t_J)$
        (so $G>0$). It reports size and power of (a) at the pool sizes used.
  \item \emph{Procedure comparison.} Compare \sfsp{} and \dfsp{} as prediction procedures across $n$,
        on $\mathcal T$. The crossover budget is reported as an empirical finding, not predicted from
        upper bounds.
\end{enumerate}
A failure to reject in (a) is evidence of low power or small $G$. It is not evidence that $G=0$.

\subsection{E3: Label acquisition}

Compare random reveal, the baselines of Section~\ref{sec:acquisition}, and the plug-in rule from
\eqref{eq:design}. All designs select from the pool only and are scored on the same test set
$\mathcal T$, so no design can move hard comparisons out of its own evaluation. Diagnostics: the
estimated maps $\hat\sigma(z)$, $\hat\gamma_1(z)$ and $\hat q(z)$, and where labels were placed.

\subsection{E4: Stress tests of the no-harm property}

Two kinds of corruption are kept separate.
\begin{itemize}
  \item \emph{Offset corruption:} corrupt only $\mJ$ and keep $Z$ fixed. (a) Constant logit shift
        $b$; (b) overconfidence, $\operatorname{logit}\mJ\mapsto\tau\operatorname{logit}\mJ$ with
        $\tau>1$; (c) region-wise reversal $\mJ\mapsto1-\mJ$ within one task stratum; (d) noise
        mixing toward a random prediction with weight $\epsilon$.
  \item \emph{Representation corruption:} apply the same corruptions to the full $\QJ$, so that $Z$
        also changes.
\end{itemize}
\textbf{Predictions.} Under offset corruption, (a) and (b) give smooth discrepancies, so \dfsp{}
should keep most of its gain. (c) gives a discontinuous discrepancy within its stratum, so
$\hat\theta_1$ should shrink there. In (d), as $\epsilon\to1$, \dfsp{} should approach the human-only
estimator on the same $Z$, as Target~\ref{tr:oracle} predicts. Under representation corruption,
no-harm relative to an uncorrupted $Z$ is \emph{not} expected; the size of this gap is reported. We
report the path of $\hat\theta$ throughout.

\subsection{E5: Judge budget (exploratory)}

With protocol J2, vary $M\in\{1,2,4,\dots,64\}$ and compare the saturation of the \dfsp{} test error
with the oracle information-loss bound of Proposition~\ref{prop:judge-budget}. No estimator-level
claim is made.

\subsection{E6: Post-training}

\textbf{Policy.} A small open-weight instruction model trained with LoRA or QLoRA
\citep{hu2022lora,dettmers2023qlora}. \todo{model, $\beta$, clipping level $D$, compute budget}. The
pipeline is
\[
  \pi_0\to\text{judge pool}\to\QJ\to\text{\dfsp{} ($n$ labels)}\to\widehat Q_H
  \to\text{pref.\ optimization}\to\pi_{\mathrm{post}}.
\]
$\pi_{\mathrm{post}}$ is a new checkpoint obtained by gradient updates.

\textbf{Arms} (all trained on the pool). DPO on the $n$ human labels only; hard RLAIF (DPO on the
judge's argmax over all $N$); soft RLAIF (soft DPO with targets $p_J$); GAPO with targets $p_J$; soft
DPO with linear-map targets \citep{vandenburg2025aligning}; soft DPO with \sfsp{} targets; \fsppo{}
(\dfsp{} targets); \fsppo{} targets inside GAPO (outside our guarantee); a hybrid with human labels
on $\cI_n$ and judge labels elsewhere, in the spirit of RLTHF; and, as an upper reference, DPO on all
$N$ pool human labels. Each arm is run at 2--3 budgets with 3 seeds.

\textbf{Policy evaluation} (on the prompts and pairs of $\mathcal T$).
\begin{itemize}
  \item[P1] \emph{Primary; the bridge quantity.} Human preference log-loss and accuracy of
        $\sigma(d_\vartheta(X))$ on test pairs. Because the loss is linear in the target, the
        observed $Y$ can replace $p_X$ and the estimate is unbiased for $R(\vartheta;p_X)$.
  \item[P2] \emph{Generation quality (proxy evaluation).} Win rate against $\pi_{\mathrm{ref}}$ on test
        prompts, scored by (i) a gold reward model trained on all pool human labels, following the
        gold-model protocol of \citet{gao2023scaling}, and (ii) an LLM judge different from the
        training judge. Both are proxies for human judgment and are reported as such.
  \item[P3] KL divergence to $\pi_{\mathrm{ref}}$ and response length (a verbosity control).
\end{itemize}
\textbf{Headline plot.} P1 and P2 versus human-label budget, with the all-human DPO reference as a
horizontal line.

\subsection{Decision gates}

\begin{description}[leftmargin=0.9cm,style=nextline]
  \item[G1] \emph{After E0/E2(a--b), on pool data only.} Is $G$ significantly positive under a test
        with verified size? If not, reframe the paper around \sfsp{}, the design theory and the
        bridge, and report the negative result.
  \item[G2] \emph{After E1.} Does \dfsp{} or \sfsp{} beat the linear map and isotonic recalibration at
        $n/N\le5\%$ on $\mathcal T$? If not, the nonparametric machinery is not paying off; center the
        paper on Targets~\ref{tr:lower}, \ref{tr:design} and~\ref{tr:e2e}.
  \item[G3] \emph{After the E6 pilot.} Does the Brier gain carry over to P1? If not, check whether the
        bounded-scorer regime of Proposition~\ref{prop:bridge} is the operative one.
\end{description}
